\subsection{例子}

\subitemhead{局部可缩空间}

定义 3. --- 称拓扑空间 B 为局部可缩的，如果 B 的每一点 b 都有一个邻域 V，使点空间 (V, b) 是可缩的（III, p. 234, 例 3）。

命题 5. --- 局部可缩的拓扑空间是可解圈的。

设 B 是一个局部可缩空间。设 b 是 B 的一点，V 是 b 的一个邻域，使 \((V, b)\) 是一个点态可缩空间。空间 V 同伦等价于一点，故 \(\pi_{1}(V, b) = \{\varepsilon_{b}\}\)（III, p. 296, 推论 3）。

为证明本命题，还需证明空间 B 局部道路连通。设 \(\sigma: V \times I \to V\) 是 \(b\) 处的一个连接像为 \(b\) 的常值映射与映射 \(Id_V\) 的点态同伦。对 \(b\) 的每个包含于 V 的邻域 W，集合 \(\sigma(W \times I)\) 是 \(b\) 的一个邻域，因为它包含 \(W = \sigma(W \times \{1\})\)；而且它道路连通，因为对每个 \((a, s) \in W \times I\)，映射 \(t \mapsto \sigma(a, ts)\) 是连接 \(b = \sigma(a, 0)\) 与 \(\sigma(a, s)\) 的一条 \(\sigma(W \times I)\) 中的道路。我们来证明诸形如 \(\sigma(W \times I)\) 的集合构成 \(b\) 的一个基本邻域系。设 \(V'\) 是 \(b\) 在 B 中的一个开邻域；于是 \(\sigma^{-1}(V')\) 是 \(\{b\} \times I\) 在 \(V \times I\) 中的一个开邻域。由于空间 I 紧，投影 \(\mathrm{pr}_1: V \times I \to V\) 是逆紧的（TG, I, p. 77, 推论 5），且 \(\mathrm{pr}_1(\complement\,\sigma^{-1}(V'))\) 是 V 的一个不包含 \(b\) 的闭子集。于是它的补 W 是 \(b\) 在 V 中的一个开邻域，使 \(\mathrm{W} \times \mathbf{I} \subset \sigma^{-1}(\mathrm{V}')\)，从而 \(\sigma(\mathrm{W} \times \mathbf{I}) \subset \mathrm{V}'\)。

数值空间或实、复 n 维射影空间的每个开子集（参见第六章与第八章）都是可解圈的，欧氏球面 \(S_{n}\)（TG, VI, p. 11, 命题 4）亦然。更一般地，每个拓扑流形（VAR R, 第二部分, p. 7, 记号与约定）都是可解圈的。事实上，这些空间都是局部可缩的。

\subitemhead{圆的庞加莱群}

拓扑空间 R 同伦于一点（III, p. 234, 例 4），因此是道路单连通的（IV, p. 340）。R 到 T = R/Z 上的典范满射 p 使之成为以 Z 为群的主覆叠（I, p. 100, 例 4）。拓扑空间 T 是可解圈的，且 (R, 0) 是 (T, 0) 的一个万有覆叠。由此我们导出群的一个典范同构 \(\pi_{1}(\mathbf{T}, 0) \to \mathbf{Z}\)：它把 0 处的圈 \(t \mapsto p(nt)\) 的类对应到 Z 的元素 n。顾及例 6（I, p. 101），我们导出下列命题。

命题 6. --- 映射 \(p \colon x \mapsto e^{2\pi ix}\)（从 \(\mathbf{R}\) 到 \(\mathbf{S}_1\) 上）使 \((\mathbf{R}, 0)\) 成为 \((\mathbf{S}_1, 1)\) 的一个万有覆叠。群 \(\pi_1(\mathbf{S}_1, 1)\) 同构于 \(\mathbf{Z}\)；圈 \(t \mapsto e^{2\pi it}\) 的类是一个生成元。对每个整数 \(n > 0\)，映射 \(z \mapsto z^n\) 使 \(\mathbf{S}_1\) 成为一个以 \(\mathbf{Z}/n\mathbf{Z}\) 为群的、\(\mathbf{S}_1\) 上的伽罗瓦覆叠。\(\mathbf{S}_1\) 的每个连通且非空的覆叠都同构于这些覆叠之一。

只需证明最后的断言。一个连通且非空的覆叠 E 在 1 上的纤维 \(E_{1}\)（E 为 \(\mathbf{S}_{1}\) 的覆叠）赋有群 \(\pi_{1}(\mathbf{S}_{1},1)\) 的一个传递作用。由于 Z 的子群都是形如 nZ 的集合（对 \(n \geqslant 0\)），可见 \(E_{1}\) 同构于与前述诸覆叠相伴的齐性集合之一。由于 \(\mathbf{S}_{1}\) 连通且局部道路连通，E 同构于这些覆叠之一。

推论. --- 映射 \(z \mapsto e^{z}\)（从 C 到 \(C^{*}\)）使 \((\mathbf{C},0)\) 成为 \((\mathbf{C}^{*},1)\) 的一个万有覆叠。群 \(\pi_{1}(\mathbf{C}^{*},1)\) 同构于 Z；圈 \(t \mapsto e^{2\pi it}\) 的类是一个生成元。对每个整数 n \textgreater{} 0，映射 \(z \mapsto z^{n}\) 使 \(C^{*}\) 成为一个以 Z/nZ 为群的、\(C^{*}\) 上的伽罗瓦覆叠。\(C^{*}\) 的每个连通且非空的覆叠都同构于这些覆叠之一。

映射 \(z \mapsto (|z|, z/|z|)\) 是空间 \(\mathbf{C}^*\) 到空间 \(\mathbf{R}_+^* \times \mathbf{S}^1\) 上的一个同胚（TG, VI, p. 10, 命题 3）；特别地，\(\mathbf{C}^*\) 是道路连通的。由于映射 \(x \mapsto e^x\)（从 \(\mathbf{R}\) 到 \(\mathbf{R}_+^*\) 上）是一个同胚，由命题 6 可知，映射 \((x, y) \mapsto e^x e^{2\pi iy}\) 使 \(\mathbf{R}^2\) 成为 \(\mathbf{C}^*\) 的一个覆叠。把 \(\mathbf{R}^2\) 与 \(\mathbf{C}\) 依映射 \((x, y) \mapsto x + iy\) 等同起来，我们断定：空间 \(\mathbf{C}\)——赋以映射 \(z \mapsto e^z\)——是 \(\mathbf{C}^*\) 的一个覆叠。由于空间 \(\mathbf{C}\) 连通且道路单连通，点覆叠 \((\mathbf{C}, 0)\) 是点空间 \((\mathbf{C}^*, 1)\) 的一个万有覆叠。

由该推论（III, p. 297）还可得出：\(\mathbf{C}^*\) 在 1 处的庞加莱群同构于 \(\mathbf{Z}\)，且类 \(\gamma\)（圈 \(t \mapsto e^{2\pi it}\) 的类）是一个生成元。

  设 n 是整数 \textgreater{} 0。映射 \(x \mapsto x^{n}\)（从 \(R_{+}^{*}\) 到其自身）是一个同胚；我们同前推出：映射 \(z \mapsto z^{n}\) 使 \(C^{*}\) 成为 \(C^{*}\) 的一个 n 次覆叠。它是以 Z/nZ 为群的伽罗瓦覆叠。最后的断言如命题 6 的证明中那样证明。

对每个整数 \(n \geqslant 0\)，环面 \((\mathbf{S}_{1})^{n}\) 是可解圈空间（IV, p. 341, 命题 1），且它在每一点的庞加莱群都同构于 \(\mathbf{Z}^{n}\)（III, p. 297, 命题 4）。

我们将在第 3 节中推广这些结果，该节专门讨论拓扑群的覆叠。

\subitemhead{实射影空间}

设 n 是整数 \textgreater{} 0。空间 \(S_{n}\) 与 \(\mathbf{P}_{n}(\mathbf{R})\) 都是连通的、非空的，且典范映射（TG, VI, p. 13）\(\varphi\colon\mathbf{S}_{n}\to\mathbf{P}_{n}(\mathbf{R})\) 使 \(S_{n}\) 成为 \(\mathbf{P}_{n}(\mathbf{R})\) 的一个主覆叠，其群为 \(\{+1,-1\}\)（I, p. 99, 例 1）。对 \(n\geqslant2\)，球面 \(S_{n}\) 是单连通的（I, p. 127, 例 3）且是可解圈的，因而是道路单连通的（IV, p. 344, 推论 1）。对 n=1，映射 \(z\mapsto z^{2}\)（从 \(S_{1}\) 到 \(S_{1}\)）在 \(S_{1}\) 中定义了一个等价关系，其诸等价类是 \(S_{1}\) 的对径点偶。由过渡到商所定义的映射 \(\varphi\colon\mathbf{S}_{1}\to\mathbf{P}_{1}(\mathbf{R})\) 定义了一个连续双射 \(\psi\colon\mathbf{S}_{1}\to\mathbf{P}_{1}(\mathbf{R})\)，使 \(\psi(z^{2})=\varphi(z)\)。由于 \(S_{1}\) 是紧空间，映射 \(\psi\) 是一个同胚（TG, I, p. 63, 推论 2）。

命题 7. --- 映射 \(x \mapsto \varphi(e^{2\pi ix})\)（从 \(\mathbf{R}\) 到 \(\mathbf{P}_{1}(\mathbf{R})\) 上）使 \((\mathbf{R},0)\) 成为 \((\mathbf{P}_{1}(\mathbf{R}),\varphi(1))\) 的一个万有覆叠。设 c: \(\mathbf{I} \to \mathbf{S}_{1}\) 是道路 \(t \mapsto e^{\pi it}\)；\(\varphi(c)\) 的类是群 \(\pi_1(\mathbf{P}_{1}(\mathbf{R}), \varphi(1))\) 的一个生成元，该群同构于 \(\mathbf{Z}\)。

对每个整数 \(n \geqslant 2\) 及 \(S_{n}\) 的每一点 x，典范映射 \(\varphi: \mathbf{S}_{n} \to \mathbf{P}_{n}(\mathbf{R})\) 使 \((\mathbf{S}_{n}, x)\) 成为 \((\mathbf{P}_{n}(\mathbf{R}), \varphi(x))\) 的一个万有覆叠。对 \(S_{n}\) 中连接 x 与 -x 的每条道路 c，\(\varphi \circ c\) 的类生成群 \(\pi_{1}(\mathbf{P}_{n}(\mathbf{R}), \varphi(x))\)，该群同构于 Z/2Z。

\subitemhead{空间关于离散群逆紧作用的商}

引理 1. --- 设 X 是一个拓扑空间，G 是一个逆紧地作用于 X 上的离散群，p 是从 X 到 X/G 上的典范投影。

设 \(x\) 是 \(X\) 的一点，\(K_x\) 是它在 \(G\) 中的固定子。设 \(U_1\) 是 \(x\) 在 \(X\) 中的一个邻域；存在一个邻域 \(U\)（\(x\) 在 \(X\) 中），它在 \(K_x\) 下稳定且包含于 \(U_1\)，使 \(g \cdot U \cap U = \varnothing\)（对 \(g \notin K_x\)），且使典范映射 \(p\) 诱导 \(U / K_x\) 到邻域 \(V\)（\(p(x)\) 在 \(X / G\) 中的邻域）上的一个同胚。

由 TG, III, p. 32 的命题 8，存在一个邻域 \(U_{2}\)（x 在 X 中），它在 \(K_{x}\) 下稳定，使得 \(g \cdot U_{2} \cap U_{2} = \varnothing\)（对 \(g \notin K_{x}\)），且使得典范映射 p 诱导一个从 \(U_{2}/K_{x}\) 到 X/G 中 \(p(x)\) 的某个邻域上的同胚。

由于典范映射 \(U_{2} \rightarrow U_{2}/K_{x}\) 是闭的（TG, III, p. 28, 命题 2），存在 x 的一个开邻域 U，它包含于 \(U_{1} \cap U_{2}\) 且在 \(K_{x}\) 下稳定。\(U_{2}\) 中由 \(K_{x}\) 定义的等价关系还是开的（TG, III, p. 10, 引理 2）。于是由 TG, I, p. 32, 命题 4 推出，典范映射诱导一个从 \(U/K_{x}\) 到 X/G 中 \(p(x)\) 的某个开邻域上的同胚，引理得证。

命题 8. --- 设 X 是一个拓扑空间，G 是一个逆紧地作用于 X 上的离散群。用 p 表示从 X 到 X/G 的典范投影。

\begin{enumerate}
\def\labelenumi{\alph{enumi})}
\item
  假设 X 道路连通，且群 G 由 X 的各点的稳定子生成。那么典范同态 \(\pi_{1}(X,x)\to\pi_{1}(X/G,p(x))\) 对 X 的每一点 x 都是满射的。特别地，若 X 道路单连通，则 X/G 亦然。
\item
  若空间 X 是可解圈的，则空间 X/G 是可解圈的。
\item
  设 N 是由满足下列条件的元素 \(g \in G\) 组成的集合：存在一条道路 \(c_{g} \colon I \to X\)，使 \(c_{g}(0) = x\)、\(c_{g}(1) = g \cdot x\)，且使 X/G 中的圈 \(p \circ c_{g}\) 严格同伦于一条常值圈；N 是 G 的一个子群。设 \(g \in G\)，且存在 X 的一点 y 使 \(g \cdot y = y\)；取一条连接 x 与 y 的道路 c: I → X；于是道路 \(c' = c * (g \cdot c)\) 满足 \(c'(0) = x\)、\(c'(1) = g \cdot x\) 且 \([p \circ c'] = [p \circ c][p \circ (g \cdot c)]^{-1} = e_{p(x)}\)，因此 \(p \circ c'\) 严格同伦于一条常值圈。由于 G 由 X 的各点的固定子生成，由此得出 N = G。
\end{enumerate}

设 \(c\) 是 X/G 中 \(p(x)\) 处的一个圈。由 III, p. 287 的定理 4，存在一条道路 \(\widetilde{c} \colon \mathbf{I} \to \mathrm{X}\)，它提升 \(c\) 且使 \(\widetilde{c}(0) = x\)。由于 \(p(\widetilde{c}(1)) = c(1) = p(x)\)，存在 \(g \in \mathrm{G}\) 使 \(\widetilde{c}(1) = g \cdot x\)。选取一条道路 \(c_g \colon \mathbf{I} \to \mathrm{X}\)，它连接 \(x\) 与 \(g \cdot x\)，且使 \(p \circ c_g\) 严格同伦于一条常值圈。于是道路 \(c' = \widetilde{c} * \overline{c_g}\) 是 X 中 \(x\) 处的一个圈，且 \([p \circ c'] = [c]\)。这表明同态 \(\pi_1(\mathrm{X}, x) \to \pi_1(\mathrm{X/G}, p(x))\) 是满射的。另一断言立即得证。

现在来证明断言 \(b\)。空间 X/G 是局部道路连通的（III, p. 261, 命题 8）。

设 \(x\) 是 X 的一点，\(K_x\) 是它在 G 中的稳定子。设 \(U_1\) 是 \(x\) 的一个开邻域，使典范同态 \(\pi_1(U_1, x) \to \pi_1(X, x)\) 的像退缩为单位元。由引理 1（IV, p. 349），存在 X 中 \(x\) 的一个邻域 U，它包含于 \(U_1\)、在 \(K_x\) 下稳定，使 \(g \cdot U \cap U = \varnothing\)（对 \(g \notin K_x\)），且使典范映射 \(p\) 诱导 \(U/K_x\) 到 X/G 中 \(p(x)\) 的邻域 V 上的一个同胚。

设 \(c\) 是 V 中 \(p(x)\) 处的一个圈；把定理 4（III, p. 287）应用于拓扑空间 U 与群 \(K_x\)，存在一条道路 \(\widetilde{c} \colon \mathbf{I} \to \mathbf{U}\)，它使 \(\widetilde{c}(0) = x\) 且提升 \(c\)。必然有 \(\widetilde{c}(1) = x\)，因而 \(\widetilde{c}\) 是 U 中 \(x\) 处的一个圈。依假设，\(\widetilde{c}\) 在 X 中严格同伦于一条常值圈；于是 \(c\) 亦然，且典范同态 \(\pi_1(V, p(x)) \to \pi_1(X/G, p(x))\) 平凡。这表明：若 X 是可解圈的，则 X/G 是可解圈的。

\begin{enumerate}
\def\labelenumi{\arabic{enumi})}
\setcounter{enumi}{4}
\tightlist
\item
  在欧氏平面 \(R^{2}\) 中，设 P 为以 \((1/n,0)\) 为中心且过原点的诸圆（对 \(n \in N^{*}\)）之并构成的拓扑空间。空间 P 是紧的、连通的且局部道路连通的，但不是可解圈的。群 \(\pi_{1}(P,0)\) 赋以容许拓扑后是豪斯多夫的且非离散（III, p. 337, 习题 7）。
\end{enumerate}

注 1. --- III, p. 287 的定理 4 与命题 8（IV, p. 349）在更一般的假设下仍然成立：G 是 R 上的有限维李群且逆紧地作用于 X 上。更多细节参见 D. MONTGOMERY 与 C. T. YANG，“The existence of a slice”，Annals of mathematics 65 (1957), p. 108--116；R. PALAIS，“On the existence of slices for actions of non-compact Lie groups”，Annals of mathematics 73 (1961), p. 295--323；以及 G. BREDON，Introduction to compact transformation groups，Academic Press, 1972。

